\Needspace{8\baselineskip}
\SourceLeaf{94}%
\noindent A short rearrangement of the defining equation for $\lambda'$\hfill\\
first gives \mbox{$\bigl(\frac{\lambda'}{m'}\bigr)^{2}=c^{2}$}. Here we must set

\[
\frac{\lambda'}{m'}=-c,\quad\text{but then}\quad
\frac{\lambda''}{m''}=c',\qquad \frac{\lambda'''}{m'''}=c'',\ \ldots
\]

\noindent to obtain the formulas of the present note. Consequently, we have the
equations.

\[
\nu=-u,\qquad P=\frac{(1+t)^{2}}{4t}\,\frac{1-u}{a^{(\infty)}},\qquad
Q=\frac{(1-t)^{2}}{4t}\,\frac{1+u}{a^{(\infty)}},
\]

\noindent as follows from the final formulas in Art.\,19 in\,the\,cited work.\ By adding and
subtracting the last two formulas, we find

\[
A^{(0)}=\frac{h-u}{a^{(\infty)}},\qquad A^{(1)}=\frac{1-hu}{a^{(\infty)}}.
\]

Second, the problem solved here of expanding
$(h-\cos\varphi)^{-\frac{3}{2}}$ falls under the general problem treated in the Dissertatiuncula (p.\,35 ff.).
No.\,7 there gives \textsc{Euler}'s three-term
relations, by which one can calculate from the first two coefficients
$A^{(0)}$, $A^{(1)}$ all the following coefficients in succession, $A^{(2)}$,
$A^{(3)}$, $\ldots$.\ For the quotients $B^{(k)}$ of the
$A$'s, \textsc{Euler}'s recurrence formulas yield the formulas established here by \textsc{Gauss}
relating $B^{(k)}$ and $B^{(k+1)}$.

\Needspace{20\baselineskip}
Finally, it is worth noting that the continued fraction given for $B^{(k)}$
is a direct case of the general formula [25] in Art.\,12 of the »Disquisitiones
generales circa seriem infinitam etc.« (Collected Works III p.\,134), in which
the quotient of two hypergeometric functions is expanded into a continued fraction.
In the notation explained there, the formula for $B^{(k)}$
would read as follows.

\[
B^{(k)}=\frac{k+\frac{3}{2}}{k+1}\,t\mathord{.} G\bigl(-\tfrac{1}{2},k-\tfrac{1}{2},k+1;t^{2}\bigr)
\]

\noindent or equivalently

\[
B^{(k)}=\frac{k+\frac{3}{2}}{k+1}\,t\mathord{.}
\frac{F\bigl(-\tfrac{1}{2},k+\tfrac{1}{2},k+2;t^{2}\bigr)}{F\bigl(-\tfrac{1}{2},k-\tfrac{1}{2},k+1;t^{2}\bigr)}.
\]

\begin{flushright}
\textsc{Fricke.}
\end{flushright}
